\bookchapter{Глава 43: Инженерия диалога — техника эффективного общения с AI}

\chapterintro{Диалог — основной способ взаимодействия в AI-программировании. Но имейте в виду: «диалог» — это не «болтовня». Эффективный диалог в AI-программировании — структурированная техника коммуникации, которую мы называем «инженерией диалога». Она изучает, как за минимальное число реплик получить код максимально высокого качества. Если прямо, вся эта глава учит вас «ясно говорить» с AI.}

\booksection{Многослойная стратегия промптов}

Диалог с AI можно разделить на три уровня, и каждый выполняет свою функцию. Поняв эти три уровня, вы сможете собирать эффективные промпты, как конструктор из кубиков.

\begin{center}
\begin{tabularx}{\linewidth}{|X|X|X|X|}
\hline
\textbf{Уровень} & \textbf{Функция} & \textbf{Пример} & \textbf{Частота} \tabularnewline
\hline
Системный уровень & Задаёт роль AI и глобальные ограничения & «Ты — эксперт по React, код должен быть на TypeScript» & Задаётся один раз в начале проекта \tabularnewline
Уровень задачи & Описывает конкретную текущую задачу & «Сделай компонент формы с полями имени пользователя и пароля» & Даётся в начале каждой задачи \tabularnewline
Уровень обратной связи & Оценивает и корректирует результат AI & «В поле пароля нет переключателя показать/скрыть, добавь его» & Даётся по качеству результата \tabularnewline
\hline
\end{tabularx}
\end{center}

\begin{center}\small\color{gray}Три уровня диалоговых промптов\end{center}

Промпт системного уровня обычно пишут в файле памяти проекта (например, CLAUDE.md) и не повторяют в каждом диалоге. Уровень задачи вы используете чаще всего — он должен содержать три элемента: «что сделать», «какими технологиями», «какие ограничения». Уровень обратной связи требует больше всего мастерства: хорошая обратная связь позволяет AI внести точные правки, а плохая лишь заставляет его угадывать заново.

\begin{tcbinfo}{Совет: три приёма обратной связи}
Практичный шаблон обратной связи — «указать проблему + объяснить ожидание + дать контекст». Например: «Обработка ошибок в строке 25 ограничивается console.log (указание проблемы); для продакшена нужен единый обработчик ошибок и понятные пользователю сообщения (ожидание); в проекте есть готовая функция handleError (контекст)». Такая обратная связь в десять раз эффективнее, чем «поправь обработку ошибок».

\end{tcbinfo}

\booksection{Диалог при отладке: как описать ошибку, чтобы AI быстро её нашёл}

Отладка — один из самых частых сценариев в AI-программировании. Когда вы наткнулись на баг, то, как вы опишете проблему AI, напрямую определяет, сможет ли он быстро найти и починить её. Вот четыре ключевых момента диалога при отладке:

\begin{enumerate}
  \item Покажите AI «сырую» информацию, а не свой пересказ. Прикрепите к сообщению журналы ошибок, трассировки исключений и соответствующий код через @ — вместо описаний вроде «там какая-то ошибка выскочила». Сырая информация даёт AI больше точности, чем ваша интерпретация.
  \item Опишите разницу между «ожидаемым поведением» и «фактическим». «Я ожидаю, что после нажатия кнопки произойдёт переход на главную страницу (ожидание), но фактически после нажатия страница никак не реагирует (факт)» — такое описание намного полезнее, чем «кнопка сломалась».
  \item Дайте шаги воспроизведения. «Сначала ввожу имя пользователя, затем нажимаю кнопку входа — и появляется эта ошибка». Шаги воспроизведения позволяют AI смоделировать ваш путь действий и легче найти проблему.
  \item Скажите, что вы уже пробовали. «Я уже проверил сетевой запрос: возвращается 200, данные в порядке» — это помогает AI отсечь уже проверенные направления и не тратить время на повторные проверки.
\end{enumerate}

Плохое описание бага

\begin{lstlisting}
Кнопка не реагирует на нажатие, посмотри
\end{lstlisting}

Хорошее описание бага

\begin{lstlisting}
После нажатия кнопки входа страница не реагирует. Ожидаемое поведение: после нажатия — переход на страницу /home Фактическое поведение: после нажатия никакой реакции, в консоли ошибок нет Шаги воспроизведения: 1. На странице входа ввести имя пользователя admin и пароль 123456 2. Нажать кнопку «Вход» 3. Страница не реагирует Что уже проверено: - Сетевой запрос проходит нормально, возвращается 200 - Событие onClick привязано - Проверен код компонента LoginButton Связанные файлы: @src/components/LoginButton.tsx @src/pages/Login.tsx
\end{lstlisting}

\booksection{Стратегии итераций: маленькие шаги против большого рефакторинга}

В AI-программировании есть две базовые стратегии итераций — для разных сценариев. Понимание их различий поможет правильно выбирать в разных ситуациях.

Стратегия «маленькие шаги» подходит для разработки функций и исправления багов. Её суть: каждый раз просить AI сделать одно небольшое изменение и переходить к следующему только после проверки. Плюс: каждый шаг проверяем и откатываем, риск под контролем. Минус: скорость ниже, нужны частые диалоги. Подходит для: разработки новых функций, поиска оптимального решения в условиях неопределённости, осторожных итераций, когда вы не вполне уверены в результате AI.

Стратегия «большой рефакторинг» подходит для масштабных изменений с понятной целью. Её суть: один раз описать полную цель и позволить AI самостоятельно спланировать и выполнить всё изменение. Плюс: скорость — AI видит общую картину и принимает более удачные архитектурные решения. Минус: высокий риск — если AI понял неверно, исправление обойдётся дорого. Подходит для: чётких обновлений архитектуры, крупной миграции кода, рефакторинга с надёжным покрытием тестами.

\begin{center}
\begin{tabularx}{\linewidth}{|X|X|X|}
\hline
\textbf{Критерий} & \textbf{Маленькие шаги} & \textbf{Большой рефакторинг} \tabularnewline
\hline
Объём одного изменения & Одна функция или один баг & Один модуль или один архитектурный слой \tabularnewline
Частота проверки & Сразу после каждого изменения & Единая проверка после всех изменений \tabularnewline
Число реплик диалога & Много (5–10 на функцию) & Мало (1–3 на модуль) \tabularnewline
Риск & Низкий, каждый шаг можно откатить & Высокий, исправление отклонений дорого \tabularnewline
Подходит для & Новые функции, поисковая разработка & Рефакторинг с ясной целью, миграция \tabularnewline
Предпосылки & Особых требований нет & Нужно покрытие тестами \tabularnewline
\hline
\end{tabularx}
\end{center}

\begin{center}\small\color{gray}Сравнение двух стратегий итераций\end{center}

На практике обе стратегии часто комбинируют. Например, при рефакторинге модуля сначала с помощью «большого рефакторинга» попросите AI предложить общий план, проверьте его, а затем выполняйте по шагам в режиме «маленьких шагов». Так вы сохраняете общую перспективу и контролируете риск исполнения.

\booksection{Управление контекстом в многоходовом диалоге}

По мере роста числа реплик нужно активно управлять контекстом диалога, чтобы избежать проблемы «AI забыл раннюю информацию». Вот несколько практичных приёмов управления контекстом.

Первое — «заякорите» ключевую информацию. В начале диалога ясно зафиксируйте ключевые ограничения и вовремя повторяйте их дальше. Например: «Помни: наши данные хранятся в localStorage, база данных недоступна» — повторяйте это ограничение в соответствующих запросах на изменения, чтобы AI не забыл.

Второе — используйте приём «итог — продолжение». Когда диалог становится слишком длинным, попросите AI «подвести итог текущего прогресса и списка задач», затем начните новый диалог с этого итога. Так вы сохраняете ключевой контекст и избегаете деградации длинного диалога.

В-третьих, «одна беседа — одна задача». Это самый простой и одновременно самый действенный способ управлять контекстом. В каждой беседе работайте только над одной задачей или одним функциональным узлом, а завершив — открывайте новую беседу для следующего. Это естественным образом ограничивает длину диалога и сохраняет контекст ясным.

В-четвёртых, используйте встроенные функции управления контекстом. В большинстве инструментов для программирования с ИИ есть кнопки «Очистить контекст» или «Новая сессия» — когда чувствуете, что качество ответов ИИ падает, не раздумывайте: открывайте новую сессию. Потерять часть истории беседы гораздо лучше, чем продолжать бороться с деградировавшим контекстом.

\begin{tcbpractice}{Попробуйте сами: побудьте «пациентом» для своего ИИ}
Специально сверстайте в браузере маленькую страницу и заложите в неё баг (например, кнопка не реагирует). Затем, следуя формату «хорошего описания бага» из этой главы, один раз опишите проблему ИИ полностью: исходные данные, ожидаемое и фактическое поведение, шаги воспроизведения, что уже проверено. Ощутите разницу в скорости, с какой ИИ находит проблему, — и вы поймёте, почему умение описывать проблему само по себе — мастерство.

\end{tcbpractice}

\begin{tcbsummary}{Краткие итоги главы}
\begin{itemize}
  \item Три слоя промпта: системный (пишется в файл памяти), слой задачи (что делаем + технологии + ограничения), слой обратной связи (указать проблему + ожидания + контекст)
  \item Четыре опоры отладки: исходные данные, ожидаемое vs фактическое, шаги воспроизведения, что уже проверено
  \item Мелкие шаги и быстрый ритм: низкий риск, много итераций — подходит для новых функций; крупный рефакторинг: высокий риск, мало итераций — требует опоры на тесты
  \item Четыре приёма управления контекстом: зафиксировать ключевую информацию, «резюме и продолжение», одна беседа — одна задача, решительно открывать новую сессию
\end{itemize}

\end{tcbsummary}
